% !TeX root = ../../Book1.tex
% 纲要标题：### 12.1 一元多项式
% 纲要位置：计划纲要.md，第 444 行。

\section{一元多项式}
\label{b1:sec:1201}

% 纲要内容（仅作写作计划，尚非教材正文）：
%
% * 域上的带余除法；整环上的因式定理与根数估计；根与重数
% * 形式导数与重根判别
% * Lagrange 插值与中国剩余定理
% * 代数闭域中的分解与代数基本定理
% * 代数基本定理在此先行引用；利用实闭域刻画的证明归附录 E.2，须先有第19章的有限 Galois 对应及第一编的 Sylow 理论
%

同一个多项式可以被看作系数的有限序列，也可以在给定元素处求值。前一种观点适合讨论整除，后一种观点适合讨论方程的根。带余除法把这两种观点联系起来：除以 \(x-a\) 的余数，恰好就是在 \(a\) 处的值。本节的系数域记为 \(K\)；涉及更大的域时，会明确说明系数与根分别位于哪里。

\subsection{带余除法、根与重数}
\label{b1:subsec:1201:01}

\cref{b1:prop:field-poly-division}已经证明：若 \(g\in K[x]\) 非零，则任意 \(f\in K[x]\) 唯一写成 \(f=qg+r\)，其中 \(r=0\) 或 \(\deg r<\deg g\)。这里能够不断消去最高次项，是因为 \(g\) 的首项系数在域中可逆。以下先把这一结论用于一次多项式。

\begin{definition}\label{b1:def:polynomial-root}
设 \(K\) 为域，\(f\in K[x]\)，\(a\in K\)。若代入所得的元素 \(f(a)\) 为零，就称 \(a\) 为 \(f\) 的一个\term[gen]{根}[root]。若 \(f\ne0\)，且对某个正整数 \(m\) 和某个 \(g\in K[x]\) 有
\[
f=(x-a)^m g,\qquad g(a)\ne0,
\]
就称 \(m\) 为根 \(a\) 的\term[chongshu]{重数}[multiplicity]。重数为一的根称为\term[dangen]{单根}[simple root]，重数至少为二的根称为\term[chonggen]{重根}[multiple root]。
\end{definition}

若 \(L/K\) 为域扩张，我们将 \(f\) 视为 \(L[x]\) 中的多项式，并用同一定义谈 \(L\) 中的根及其重数。把系数域换成整环 \(R\) 时，也可用上述条件定义 \(R\) 中的根及其重数；下面的因式定理与根数估计在这个范围内仍然成立。

\begin{proposition}[因式定理与根数估计]\label{b1:prop:factor-root}
设 \(R\) 为整环，\(f\in R[x]\) 为非零多项式，\(a\in R\)。则 \(f(a)=0\) 当且仅当 \(x-a\mid f\)。多项式 \(f\) 的每个根都有唯一重数；若 \(\deg f=n\)，则 \(f\) 在 \(R\) 中的根按重数计至多有 \(n\) 个。
\end{proposition}
\begin{proof}
这里的除数 \(x-a\) 首一，不需要对 \(R\) 中的非零系数作除法。对每个 \(i\ge1\)，恒等式
\[
x^i-a^i=(x-a)\sum_{j=0}^{i-1}x^{i-1-j}a^j
\]
说明 \(f-f(a)\) 是 \(x-a\) 的倍数，故 \(f=(x-a)q+f(a)\)。反过来，\(x-a\) 的任何倍数在 \(a\) 处都为零，这就证明了因式定理。若 \(a\) 是根，就可以反复除去 \(x-a\)。每次除法都使非零商的次数减少一，所以至多进行 \(\deg f\) 次。记除去的次数为 \(m\)，则最后得到 \(f=(x-a)^m g\)、\(g(a)\ne0\)。若又有 \(f=(x-a)^r h\)、\(h(a)\ne0\)，设 \(m<r\)，在整环 \(R[x]\) 中消去 \((x-a)^m\)，便有 \(g=(x-a)^{r-m}h\)，与 \(g(a)\ne0\) 矛盾。交换两种分解也排除 \(r<m\)，故重数唯一。

若 \(a\ne b\)，从 \(f=(x-a)^m g\) 可见，\(b\) 是 \(f\) 的根当且仅当它是 \(g\) 的根，因为 \((b-a)^m\ne0\)，且整环中非零元素相乘仍非零。再将 \(g\) 中所有 \(x-b\) 因子除去，所得最后因子在 \(b\) 处非零，乘以 \((x-a)^m\) 后仍然非零，因此 \(b\) 的重数也不变。对任意有限个不同的根，逐一除去它们对应的因子，次数总共下降各重数之和，且不会降为负数，故这一和至多为 \(n\)。因此不同根本身也不可能超过 \(n\) 个，根数估计随之成立。
\end{proof}

例如在 \(\mathbb Q[x]\) 中，\((x-1)^3(x+2)\) 有两个不同的根，但按重数共有四个根。根数估计要求多项式非零，也要求系数处于整环；它不要求所有非零系数可逆，所以也适用于 \(\mathbb Z[x]\)。在 \(\mathbb Z/8\mathbb Z\) 中，\(x^2-1\) 在 \(\bar1,\bar3,\bar5,\bar7\) 处都为零，二次多项式却有四个不同的根。失败来自零因子：不同根之差虽非零，却可能是零因子，因而从 \((b-a)^m g(b)=0\) 不能推出 \(g(b)=0\)。

\subsection{形式导数与重根判别}
\label{b1:subsec:1201:02}

导数在这里由系数定义，不借助极限。它在任意特征下都有意义，但正特征会改变整数系数的取值。

\begin{definition}\label{b1:def:formal-derivative}
设 \(K\) 为域，\(n\geq0\) 为整数，\(a_0,\ldots,a_n\in K\)，并令 \(f=\sum_{i=0}^n a_i x^i\)。定义 \(f\) 的\term[xingshi-daoshu]{形式导数}[formal derivative]为
\[
f'\coloneqq\sum_{i=1}^n i a_i x^{i-1}.
\]
这里 \(i a_i\) 是 \(a_i\) 的 \(i\) 次相加，整数系数在 \(K\) 中解释。
\end{definition}
\symidx{\(f'\)：多项式的形式导数}

形式导数记录的是平移变量后的一次项。设 \(L/K\) 为域扩张，\(a\in L\)，\(T\) 为不定元。按二项式公式展开 \(f(a+T)\)，所得常数项为 \(f(a)\)，一次项为 \(f'(a)T\)，所以存在 \(h\in L[T]\)，使
\[
f(a+T)=f(a)+f'(a)T+T^2h(T).
\]
若 \(f\ne0\) 且 \(f(a)=0\)，常数项已经消失；再要求 \(f'(a)=0\)，便消去一次项，恰好得到 \(T^2\mid f(a+T)\)，也就是 \((x-a)^2\mid f(x)\)。因此 \(a\) 是重根恰好要求这两项同时消失。下面的因式证明是这一现象的另一种表述。这个解释只用多项式恒等式。在特征 \(p>0\) 时，\((a+T)^p=a^p+T^p\)，一次项确实不存在，这便是 \((x^p)'=0\) 的代数意义。

\ProseParagraphBreak
定义直接给出 \((f+g)'=f'+g'\) 及 \((cf)'=cf'\)。对单项式 \(ax^i,bx^j\)，由系数恒等式 \((i+j)ab=iab+jab\) 得到
\[
\bigl(ax^i\cdot bx^j\bigr)'=(ax^i)'bx^j+ax^i(bx^j)'.
\]
对有限项求和，便得到乘积法则 \((fg)'=f'g+fg'\)。因而也有 \(\bigl((x-a)^m\bigr)'=m(x-a)^{m-1}\)。

\begin{proposition}[重根判别]\label{b1:prop:multiple-root}
设 \(K\) 为域，\(f\in K[x]\) 非零，\(L\) 是包含 \(K\) 的域。元素 \(a\in L\) 是 \(f\) 的重根，当且仅当 \(f(a)=f'(a)=0\)。此外，以下条件等价：
\begin{enumerate}
\item \(\gcd(f,f')=1\)，其中最大公因式取首一形式；
\item 在任何包含 \(K\) 的域中，\(f\) 都没有重根。
\end{enumerate}
\end{proposition}
\begin{proof}
若 \(f(a)=0\)，在 \(L[x]\) 中由\cref{b1:prop:factor-root}写成 \(f=(x-a)g\)。求导后代入 \(a\)，得到 \(f'(a)=g(a)\)。因此 \(f'(a)=0\) 等价于 \(x-a\mid g\)，也就等价于 \((x-a)^2\mid f\)。

若 \(\gcd(f,f')=1\)，由于 \(K[x]\) 是主理想整环，\cref{b1:prop:pid-bezout}给出 \(u,v\in K[x]\)，使 \(uf+vf'=1\)。这个等式在 \(L[x]\) 中仍成立，所以不可能有 \(a\in L\) 同时使 \(f(a)\) 和 \(f'(a)\) 为零。

反过来，若最大公因式有正次数，在其中选一个不可约因子 \(q\)。\cref{b1:thm:euclidean-pid}及\cref{b1:prop:pid-bezout}说明 \(K[x]/(q)\) 是域：若 \(q\nmid h\)，则 \(\gcd(q,h)=1\)，某个 Bézout 等式模 \(q\) 后给出 \(h+(q)\) 的逆。常数域 \(K\) 嵌入这个商域，因为非零常数不可能被正次数的 \(q\) 整除。令 \(a=x+(q)\)，便有 \(q(a)=0\)；这个商域的作用是向 \(K\) 中形式地添入 \(q\) 的一个根，相关构造将在\secref{b1:sec:1401}系统研究。由于 \(q\mid f\) 及 \(q\mid f'\)，又有 \(f(a)=f'(a)=0\)，故 \(a\) 是其中的重根。
\end{proof}

这里的“任何包含 \(K\) 的域”不能改成仅检查 \(K\) 本身。例如 \((x^2+1)^2\) 在 \(\mathbb Q\) 中没有根，在 \(\mathbb C\) 中却有重根。判别条件只涉及系数域 \(K\) 中的最大公因式，却能控制所有扩域中的重根，不必先求出各个扩域中的根。

\ProseParagraphBreak
设域 \(K\) 的特征为素数 \(p\)。此时 \(f'=0\) 当且仅当 \(f\) 的所有非零系数项的指数都被 \(p\) 整除，即 \(f=\sum_j a_jx^{pj}\)：若 \(p\nmid i\)，则 \(i\cdot1_K\ne0\)，从 \(ia_i=0\) 可消去它而得 \(a_i=0\)。这是关于指数的条件，并未要求系数 \(a_j\) 具有 \(p\) 次根。因此非常数多项式也可能导数为零。例如在 \(\mathbb F_p[x]\) 中，\(x^p-1=(x-1)^p\)，它的导数为零；等式来自中间二项式系数都被 \(p\) 整除。

\ProseParagraphBreak
要进一步写成 \(f=g^p\)，还须解决系数问题：为每个 \(a_j\) 在 \(K\) 中找到 \(b_j\)，使 \(a_j=b_j^p\)。这个条件既必要也充分，因为
\[
\left(\sum_j b_jx^j\right)^p=\sum_j b_j^px^{pj}.
\]
所以 \(f'=0\) 只限制指数，系数是否有 \(p\) 次根则决定 \(f\) 能否成为 \(K[x]\) 中的 \(p\) 次幂。具体例子见\cref{b1:ex:12-primitive-linear-polynomials}。

\subsection{Lagrange 插值与中国剩余定理}
\label{b1:subsec:1201:03}

根数估计还带来另一种用途：两个次数有界的多项式，在足够多的点处相等，就必定相同。

\begin{theorem}[Lagrange 插值]\label{b1:thm:lagrange-interpolation}
设 \(K\) 为域，\(n\geq0\) 为整数，\(a_0,\ldots,a_n\in K\) 互不相同，\(b_0,\ldots,b_n\in K\)。令
\[
\ell_i(x)=\prod_{\substack{0\le j\le n\\j\ne i}}\frac{x-a_j}{a_i-a_j}
\qquad(0\le i\le n).
\]
空乘积取一。这些多项式满足 \(\ell_i(a_j)=\delta_{ij}\)，其中 \(\delta_{ij}\) 在 \(i=j\) 时为一，否则为零。存在唯一的次数不超过 \(n\) 的多项式 \(f\in K[x]\)，使 \(f(a_i)=b_i\) 对所有 \(0\leq i\leq n\) 成立。它由
\begin{equation}\label{b1:eq:lagrange-interpolation}
f(x)=\sum_{i=0}^n b_i\ell_i(x)
\end{equation}
给出；这里也允许 \(f=0\)。
\end{theorem}
\begin{proof}
由于各 \(a_i\) 互异，每个分母均非零，故 \(\ell_i\) 有意义。它在 \(a_i\) 处每个因子都等于一，在其余 \(a_j\) 处有一个因子为零，所以 \(\ell_i(a_j)=\delta_{ij}\)。每个 \(\ell_i\) 的次数至多为 \(n\)，故所列 \(f\) 有所需取值，次数也至多为 \(n\)。这就是用每次只保留一个指定值的插值基多项式，拼出全部指定值。若 \(g\) 也是这样的多项式，则 \(f-g\) 的次数至多为 \(n\)，却有 \(n+1\) 个不同的根。\cref{b1:prop:factor-root}迫使 \(f-g=0\)。
\end{proof}

这称为\term[lagrange-chazhi]{Lagrange 插值}[Lagrange interpolation]。例如在 \(\mathbb Q\) 中指定 \(f(0)=1,f(1)=2,f(2)=5\)，公式给出
\[
f(x)=\frac{(x-1)(x-2)}2-2x(x-2)+\frac{5x(x-1)}2=x^2+1.
\]
唯一性只对次数至多为二的多项式成立；加上任意 \(x(x-1)(x-2)h(x)\)，这三个值都不改变。

\ProseParagraphBreak
\phantomsection\label{b1:note:polynomial-function-distinction}
形式多项式相等，指对应系数逐项相等；它们在每个元素处的取值当然也相等，即
\[
\boxed{\ f=g\;\text{于}\;K[x]\quad\Longrightarrow\quad
f(a)=g(a)\;\text{对每个}\;a\in K\ }.
\]
反向一般不成立。设 \(p\) 为素数。在 \(\mathbb F_p[x]\) 中，\(x^p-x\) 的首项系数为一，所以不是零多项式，却在每个 \(a\in\mathbb F_p\) 处取值为零：\(a=0\) 时直接成立，\(a\ne0\) 时，\cref{b1:thm:lagrange}用于 \(p-1\) 阶群 \(\mathbb F_p^\times\) 给出 \(a^{p-1}=1\)。因此它与零多项式定义同一个函数。要由取值推出多项式相等，还须用根数估计核对次数界；例如在 \(\mathbb F_p\) 上，两个次数均至多为 \(p-1\) 的多项式若处处取值相同，其差便有 \(p\) 个不同根而次数小于 \(p\)，只能是零多项式。

\ProseParagraphBreak
这个插值公式也可以从\cref{b1:thm:crt}理解。对定理中的互异元素 \(a_0,\ldots,a_n\)，令 \(I_i=(x-a_i)\subseteq K[x]\)。若 \(i\ne j\)，则 \((x-a_i)-(x-a_j)=a_j-a_i\) 是非零常数，故 \(I_i+I_j=K[x]\)。中国剩余定理给出同构
\[
K[x]\Big/\left(\prod_{i=0}^n(x-a_i)\right)
\longrightarrow K^{n+1},\qquad
[g]\longmapsto\bigl(g(a_0),\ldots,g(a_n)\bigr).
\]
定理中的插值基多项式 \(\ell_i\) 在这个商环中对应第 \(i\) 个坐标幂等元，因为它的取值向量只有第 \(i\) 项为一，其余项全为零。中国剩余定理负责实现任意指定值，确定相应的唯一陪集；带余除法再为这个陪集选出唯一的低次代表。具体地，对 \(P=\prod_{i=0}^n(x-a_i)\) 作带余除法，所得余式的次数至多为 \(n\)。若两个这样的代表属于同一陪集，它们的差是 \(P\) 的倍数，而非零倍数的次数至少为 \(n+1\)，所以两者必相等。

\subsection{代数闭域中的分解与代数基本定理}
\label{b1:subsec:1201:04}

根数估计限制已有根的数量，却不保证根一定存在。现在换一个问题：在哪些域中，每个非常数多项式都至少有一个根？

\begin{definition}\label{b1:def:algebraically-closed}
若域 \(K\) 上每个非常数多项式在 \(K\) 中都有根，就称 \(K\) 为\term[daishu-bi-yu]{代数闭域}[algebraically closed field]。
\end{definition}

定义只要求至少存在一个根，却会自动产生完全分解：提出一个一次因子以后，商仍在 \(K[x]\) 中；只要它仍是非常数多项式，就能再次使用同一个存在条件，直到次数降为零。

\begin{proposition}\label{b1:prop:algebraically-closed-splitting}
域 \(K\) 代数闭，当且仅当每个次数为 \(n\ge1\) 的多项式 \(f\in K[x]\) 都能写成
\[
f(x)=a_n\prod_{i=1}^n(x-\alpha_i),\qquad \alpha_i\in K,
\]
其中 \(a_n\) 是 \(f\) 的首项系数。这时按重数计共有 \(n\) 个根，根的多重集合由 \(f\) 唯一确定。
\end{proposition}
\begin{proof}
若 \(K\) 代数闭，选 \(f\) 的一个根 \(\alpha_1\)。\cref{b1:prop:factor-root}给出 \(f=(x-\alpha_1)g\)，而 \(\deg g=n-1\)。对次数归纳，直到商为非零常数，得到所需分解，常数因子由首项系数确定。每个根出现的次数恰为\cref{b1:prop:factor-root}中唯一确定的重数，因此根的多重集合唯一。反过来，所列分解直接给出任一非常数多项式的一个根。
\end{proof}

对通常的复数域，根的存在还要另作论证。以下结论称为\term[daishu-jiben-dingli]{代数基本定理}[fundamental theorem of algebra]；它不是复数的定义。此处先引用这一结论；利用实分析的中间值定理和有限 Galois 理论的证明见\secref{b1:sec:E02}。

\begin{theorem}[代数基本定理]\label{b1:thm:fundamental-theorem-algebra}
复数域 \(\mathbb C\) 是代数闭域。因此，每个次数为 \(n\ge1\) 的复系数多项式都有恰好 \(n\) 个复根，按重数计，并可分解为一次因子的乘积。
\end{theorem}

不可约性讨论必须始终注明系数域；例如 \(x^2+1\) 在 \(\mathbb R\) 中无根，在 \(\mathbb C[x]\) 中却等于 \((x-i)(x+i)\)。
